
\section*{Part II: The Universal Theory of Growing from 0}
\addcontentsline{toc}{section}{Part II: The Universal Theory of Growing from 0}

\subsection{8. 0 is not the absence of assets}

``0 rupees'' is not the same as 0 resources. A technically capable founder may possess: expertise, credibility from prior work, time, a network, a laptop, access to inexpensive software, and the ability to solve an urgent problem. These are not cash, but they are capital in the economic sense.

The 1st task is to convert \textbf{scarce founder time into verified market learning and positive cash flow}, without making liabilities that could destroy the venture before it learns anything.

A useful ordering is:

\begin{enumerate}[leftmargin=*,itemsep=1.5pt]
  \item \textbf{Earn by solving a problem you understand.}
  \item \textbf{Discover a repeatable problem across customers.}
  \item \textbf{Standardize the solution where standardization does not degrade value.}
  \item \textbf{Hire and automate the repeatable components.}
  \item \textbf{Turn repeated delivery into intellectual and relational assets.}
  \item \textbf{Use those assets to enter 1 adjacent market at a time.}
  \item \textbf{Fund long-horizon capabilities from a durable base.}
\end{enumerate}

Do not mistake the order for a rule that all businesses begin with consulting. Some products can be built and distributed cheaply from day 1. The principle is \textbf{early cash and fast feedback}, not a universal requirement to offer services.

\subsection{9. Why doing things that do not scale is rational}

Paul Graham's essay \href{https://paulgraham.com/\allowbreak ds.html}{``Do Things That Don't Scale''} argues that founders often have to recruit early users manually. The deeper reason is informational. Before a company understands the customer, it does not know which steps should be automated. The founder's unusually high-touch behavior is partly a measurement instrument.

Teaching the 1st 10 learners oneself reveals where explanations fail, how motivation varies, whether projects are too difficult, what people value, and why they abandon courses. Selling the 1st 3 college partnerships oneself reveals who really approves budgets, what procurement requires, whether professors are threatened by outside trainers, and where implementation breaks.

Unscalable behavior is justified when it produces \textbf{transferable learning}. It becomes dangerous when it is merely a permanently subsidized premium service.

\textbf{\textbf{[Decision rule]}} Record what each founder-intensive interaction teaches. For every repeated task, ask whether its value comes from founder judgment, human care, or routine administration. Preserve the 1st 2; progressively systematize the 3rd.

\subsection{10. A sequence of constraints, not a sequence of titles}

At different sizes, a business has different binding constraints.

\begin{table*}[t]
\centering
\small
\caption{A Sequence of Organizational Constraints, Key Questions, and Primary Artifacts}
\label{tab:ch36_constraints_sequence}
\begin{tabularx}{\textwidth}{p{2.6cm} p{3.2cm} p{4.5cm} X}
\toprule
\textbf{ Organizational state } & \textbf{ Real constraint } & \textbf{ Management question } & \textbf{ Primary artifact } \\
\midrule
 No customers  &  no proof of demand  &  Who will commit money now?  &  paid experiment  \\
 1st successful customers  &  no reproducibility  &  Can the result be delivered again?  &  delivery playbook  \\
 Repeat demand  &  founder capacity  &  Can someone else deliver quality?  &  trained 2nd instructor / operator  \\
 Operationally repeatable  &  distribution bottleneck  &  Can customers arrive predictably?  &  channel economics  \\
 Several profitable offerings  &  managerial complexity  &  Can leaders own outcomes?  &  operating cadence and P\&Ls  \\
 Strong commercial platform  &  capital allocation  &  Which adjacency has the best strategic return?  &  investment committee process  \\
 Research ambition  &  scientific credibility  &  Can independent experts respect the work?  &  researchers, peer review, grants  \\
 Physical-degree ambition  &  financing and legitimacy  &  Can the institution sustain the campus and obey the law?  &  legal architecture, capital stack, academic plan  \\
 Mature institute  &  succession and renewal  &  Will excellence survive leadership change?  &  governance, endowment policy, cultural transmission  \\
\bottomrule
\end{tabularx}
\end{table*}

A company cannot use better marketing to solve weak pedagogy. It cannot use a campus to solve an absent research culture. Each stage requires solving the actual constraint.

\subsection{11. Product-market fit is a condition, not a congratulatory slogan}

Product-market fit is not a single threshold. There can be \textbf{offer fit}, \textbf{channel fit}, \textbf{unit-economics fit}, \textbf{operational fit}, and \textbf{institutional fit}. For Bayes, a profitable cohort with 12 highly satisfied learners may establish initial offer fit. It does not establish that 2,000 learners will arrive affordably, that non-founder faculty can deliver, or that the course creates a research institution.

Measure 5 distinct proofs:

\begin{itemize}[leftmargin=*,itemsep=1.5pt]
  \item \textbf{Demand:} strangers pay at a price that can support the model.
  \item \textbf{Outcome:} genuine learners develop demonstrable competence.
  \item \textbf{Repeatability:} later cohorts achieve comparable results.
  \item \textbf{Acquisition:} the next cohort does not rely on exhausting the founder's personal network.
  \item \textbf{Delegation:} at least part of the experience survives transfer to other instructors and staff.
\end{itemize}

The last proof is particularly important for institutional evolution. A founder who stays indispensable has built a strong personal practice, not yet an institution.

\subsection{12. Sequence experiments by reversibility}

Before making any meaningful commitment, classify it:

\begin{itemize}[leftmargin=*,itemsep=1.5pt]
  \item \textbf{Type A: Cheap and reversible.} A weekend workshop, a small paid pilot, content, visiting expert, a short enterprise discovery engagement.
  \item \textbf{Type B: Moderately reversible.} A 6-month faculty contract, a learning platform, a year-long college partnership, a small leased classroom.
  \item \textbf{Type C: Difficult to reverse.} Long leases, multi-year payroll commitments, high-interest debt, capital-intensive labs, land purchase, major regulatory submissions, reputationally risky degree claims.
  \item \textbf{Type D: Institution-defining.} University sponsor structure, academic freedoms, permanent campus location, investor rights that survive generations, research and IP governance.
\end{itemize}

Move quickly on A, carefully on B, very slowly on C and D. The ambition should be immense; the irreversible bets should be exceptionally well justified.

\subsection{13. The compounding test for an adjacent business}

A new vertical is attractive when it shares at least 2 assets with the existing business: customer relationships, subject-matter expertise, pedagogy, faculty, distribution, content/\allowbreak IP, tooling, brand, research questions, or infrastructure.

Suppose Bayes teaches applied AI to working engineers. Enterprise L\&D might share curriculum and instructors, but require distinct procurement and account management. College electives might share content, but demand faculty enablement and lengthy payment collection. A cybersecurity product may share customers but not deep founder expertise. A research software tool may share techniques but create an entirely different product/\allowbreak support burden.

A ``related'' market is not related merely because the word AI appears in both titles.

\textbf{\textbf{[Decision rule]}} Require each adjacent offer to identify its asset reuse, customer pull, missing capabilities, owner, cash requirement, downside, and expected new institutional asset. No scorecard should substitute for actual evidence of customer willingness to pay.

-


